\section{Private debt and special situations}

\begin{qfmethod}[Debt financing in private markets over the investment life cycle: leveraged loans, high-yield bonds and convertibles, and mezzanine against unitranche]\label{meth:private-debt-financing-structures}
\qffield{Statement} \emph{Over the life cycle} (Interpretation, organised from Demaria and co-authors' map of strategies \citep[p.~58, Fig.~3.2]{demaria-private-markets-2021} and Ippolito's instruments):
\begin{itemize}
\item \emph{start-up}: little debt; venture lending and leasing alongside equity, repaid from the next equity round;
\item \emph{growth}: convertible loans or preferred shares and mezzanine, whose equity features suit a company with growth but uneven cash flow;
\item \emph{buyout}: senior leveraged loans, second lien, high-yield bonds, mezzanine or a single unitranche, sized to the cash flows of a mature company;
\item \emph{maturity}: refinancing and recapitalisations when debt has been repaid;
\item \emph{distress}: rescue financing, debt-for-equity exchanges and loan-to-own (\hyperref[meth:special-situations-distressed]{special situations}).
\end{itemize}

\emph{Instruments.}
\begin{itemize}
\item \emph{Leveraged loans}: senior secured term loans arranged by banks or credit funds, an amortising Term Loan A (about six years) and bullet Term Loan B (about seven) priced over a floating rate, with a revolver for working capital; senior leverage of about 3.5--4.0 times EBITDA in normal years \citep[pp.~31--32, 192--193]{ippolito-private-capital-2020}.
\item \emph{High-yield bonds}: sub-investment-grade securities, often issued under Rule 144A, fixed-rate with non-call periods and incurrence covenants; usually at least \$150 million and structurally subordinated to operating-company debt (pp.~37--41).
\item \emph{Convertibles}: debt or preferred shares convertible into ordinary shares at a set conversion price, a bond plus a call option on the equity (pp.~57--58).
\end{itemize}

\emph{Mezzanine against unitranche} (pp.~39--45, Table~2.4):
\begin{center}\small
\begin{tabular}{@{}p{0.18\linewidth}p{0.36\linewidth}p{0.36\linewidth}@{}}
\toprule
 & Mezzanine (European form) & Unitranche \\
\midrule
Structure & separate junior tranche, second-ranking security, intercreditor agreement & one facility combining senior and junior, ranked among lenders by agreement \\
Pricing & cash plus PIK interest, often warrants & blended floating margin, slightly above a senior/junior split \\
Repayment & bullet, 6--8 years & bullet, 5--8 years \\
Prepayment & 3/2/1 percent premiums & make-whole, then 1 percent \\
Covenants & maintenance, behind senior & maintenance, comprehensive \\
For the borrower & less dilutive than equity, more flexible than bank debt & one document and lender group, certainty and speed \\
\bottomrule
\end{tabular}
\end{center}
Mezzanine lenders accept standstill periods and payment blocks in exchange for a seat in any restructuring; US mezzanine is usually unsecured and fixed-rate (pp.~41--43).

\qffield{Practice} \emph{Corran Industrial} (original): 250 of debt at 5.0 times EBITDA. A unitranche at 9 percent cash costs 22.5 of interest in year 1. Senior debt of 175 at 7 percent and mezzanine of 75 at 7 percent cash plus 5 percent PIK costs 17.5 in cash and 3.75 of PIK, an all-in 8.5 percent. The split lowers cash interest and raises year-1 cover from 2.36 to 3.03 times, at the price of two lender groups and an intercreditor negotiation. The equity returns are almost the same (2.03 times either way; \hyperref[meth:private-equity-valuation]{valuation}).

\qffield{Conditions} Terms are typical ranges from Ippolito's practitioner account, for European and US markets before 2020; they move with credit conditions.

\qffield{Pitfalls} Comparing mezzanine's cash coupon with unitranche without the PIK and warrants. Ignoring structural subordination of holding-company debt.

\qffield{See also} \hyperref[def:private-debt-features]{private debt}; \hyperref[meth:private-debt-credit-analysis]{credit analysis}.

\end{qfmethod}

\begin{qfmethod}[Analysing a private debt investment: profile, leverage, coverage, loan-to-value and covenant ratios, downside cases and the lender's decision]\label{meth:private-debt-credit-analysis}
\qffield{Statement} A private lender asks whether the borrower's cash flow can service and repay the debt through a downturn, and what it recovers if not. The profile covers the business, sponsor, purpose, ranking and security, and documentation; the ratios \citep[pp.~236--238]{ippolito-private-capital-2020}:
\begin{itemize}
\item \emph{leverage}: debt over EBITDA, generally acceptable up to about 4 times (senior), with total leverage higher in buyouts; gearing and debt to capital;
\item \emph{coverage}: interest cover (EBITDA over interest), above 2 times for senior debt and 1.5 times still acceptable; debt service cover (cash flow over interest plus scheduled principal), at least 1.5 times for less stable cash flows;
\item \emph{liquidity}: current and quick ratios;
\item \emph{asset cover}: loan to value, where repayment depends on assets.
\end{itemize}
Financial covenants should be set where a breach gives warning before the loan is impaired, and affirmative and negative covenants protect the sources of repayment (pp.~234--236).

\qffield{Practice} \emph{Corran Industrial} (original), EBITDA 50 rising to 53 in year 1; enterprise value 450:
\begin{center}\small
\begin{tabular}{@{}lrr@{}}
\toprule
 & Unitranche & Senior + mezzanine \\
\midrule
Total leverage & 5.0x & 5.0x \\
Loan to value & 55.6\% & 55.6\% \\
Year-1 interest cover (cash) & 2.36x & 3.03x \\
Debt service cover (bullet) & 1.43x & -- \\
All-in cost & 9.0\% & 8.5\% \\
Leverage if EBITDA falls 20\% & 6.25x & 6.25x \\
Interest cover if EBITDA falls 20\% & 1.88x & 2.42x \\
\bottomrule
\end{tabular}
\end{center}
\emph{Decision.} The unitranche lender is covered on interest even in the downside but breaches a 6.0 times leverage covenant, which is the point of the covenant: it forces a discussion while value remains, at 44 percent equity cushion below the loan. Lend at 5.0 times with a covenant at 6.0 times stepping down to 5.0 times, a cash sweep, and security over the shares and assets; decline 5.5 times, where a 20 percent fall leaves interest cover near 1.7 times.

\qffield{Conditions} The free cash flow is assumed to be half of EBITDA; ratio thresholds are indicative and depend on industry and stability.

\qffield{Pitfalls} Relying on enterprise value at entry multiples for recovery. Measuring leverage on adjusted EBITDA full of add-backs. Setting covenants with so much headroom that a breach comes only after impairment.

\qffield{See also} \hyperref[meth:private-debt-financing-structures]{financing structures}; \hyperref[meth:corporate-credit-analysis]{corporate credit analysis}.

\end{qfmethod}

\begin{qfmethod}[Special situations and distressed investing: characteristics and risks, distressed debt, rescue financing and restructuring options, complex situations, due diligence and recovery valuation]\label{meth:special-situations-distressed}
\qffield{Statement} \emph{Characteristics.} Special situations invest where a company or its owners face an event that public markets and normal lenders avoid: financial distress, a forced sale, a complex structure or a dislocation. Returns come from buying at a discount and resolving the situation; risks are legal, operational and timing, with heavy dependence on bankruptcy law and court practice \citep[pp.~126--128]{demaria-private-markets-2021}.

\emph{Distressed debt} \citep[pp.~127--130]{demaria-private-markets-2021}:
\begin{itemize}
\item \emph{debt-for-trading}: buying at a discount and selling after a recovery, close to hedge-fund practice and dependent on liquid debt markets;
\item \emph{debt-for-control} (loan-to-own): buying the class that will convert into the new equity, then restructuring and selling the business.
\end{itemize}
It works best under regimes like US Chapter~11, which stays payments and lets a court confirm a plan over dissenting creditors.

\emph{Voting and confirmation are different tests} \citep[secs.~1124, 1126, 1129]{us-bankruptcy-code-11}.
\begin{itemize}
\item \emph{Class acceptance} (sec.~1126(c)): a class of claims accepts when holders of at least two-thirds in amount \emph{and} more than half in number of the claims that vote accept. A holder of more than one-third of the amount voted can therefore stop its class from accepting, but holding two-thirds is not enough to make it accept if many small holders vote against.
\item \emph{Confirmation} (sec.~1129): a plan normally needs every impaired class to accept (sec.~1129(a)(8)). Otherwise the court confirms it over a dissenting impaired class only if all other conditions are met, at least one impaired class other than insiders has accepted (sec.~1129(a)(10)), and the plan does not discriminate unfairly and is fair and equitable to that class (sec.~1129(b)(1)). A class whose rights are left unaltered, or whose defaults are cured and maturity reinstated, is not impaired (sec.~1124). Each dissenting creditor must also receive at least what it would get in a chapter~7 liquidation (sec.~1129(a)(7)).
\item \emph{Fair and equitable} differs by class (sec.~1129(b)(2)). A dissenting \emph{secured} class must keep its liens and receive deferred cash payments worth at least its collateral, or have the liens attach to sale proceeds, or receive the ``indubitable equivalent''. A dissenting \emph{unsecured} class must either receive property worth its allowed claims or no junior class may receive or retain anything on account of its junior claim or interest: the absolute priority rule. Impairment alone does not decide it; what matters is what the class and those junior to it receive.
\item A secured claim is secured only up to the value of the collateral; any excess is an unsecured claim (sec.~506(a)).
\end{itemize}

\emph{Debtor-in-possession financing} (sec.~364). Ordinary-course unsecured credit is an administrative expense. If that is unavailable, the court may grant priority over other administrative expenses, a lien on unencumbered property, or a junior lien on encumbered property. A lien senior or equal to an existing lien, which \emph{primes} the pre-petition lenders, requires that credit cannot be obtained otherwise and that the existing lienholders' interest is adequately protected, with the burden of proof on the debtor (sec.~364(d)). A DIP facility therefore ranks ahead of pre-petition secured claims on their collateral only under a priming order or with their consent.

\emph{Financing alternatives in distress} \citep[pp.~291--297]{ippolito-private-capital-2020}: asset sales and sale-leasebacks; secured financing against unencumbered assets; an equity injection, often from the sponsor; debt repurchases at a discount; exchange offers, which can be coercive by stripping covenants from non-tendered bonds; out-of-court restructuring, which suffers from hold-outs; and in-court reorganisation, where debtor-in-possession financing can be given the priority described above.

\emph{Complex situations.} Non-performing loan portfolios sold by banks, orphaned or non-core assets, and companies facing a dislocation rather than insolvency (pp.~128--129). \emph{Due diligence and valuation}: the business's going-concern value under a reorganisation plan, the claims and their priority, the fulcrum security, the legal process and its timing, and the investor's ability to influence the outcome (Ippolito, pp.~292--297).

\qffield{Practice} \emph{Halden Retail} (original). Claims: a DIP facility of 40, administrative claims 10, senior secured 200, unsecured bonds 150. The waterfall rests on stated assumptions: the DIP holds a priming first lien granted under sec.~364(d) with the senior lenders' consent; the senior lenders agreed a carve-out that pays administrative claims, which a plan must otherwise pay in cash on the effective date (sec.~1129(a)(9)(A)); the senior collateral is worth at least 200, so there is no deficiency claim; and the bonds form one class.
\begin{center}\small
\begin{tabular}{@{}lrrrrr@{}}
\toprule
Reorganised value & 220 & 250 & 255 & 300 & 350 \\
\midrule
Senior secured recovery & 85.0\% & 100\% & 100\% & 100\% & 100\% \\
Unsecured recovery & 0.0\% & 0.0\% & 3.3\% & 33.3\% & 66.7\% \\
Old equity & 0 & 0 & 0 & 0 & 0 \\
\bottomrule
\end{tabular}
\end{center}
\begin{itemize}
\item \emph{Threshold}: the claims ahead of the bonds total $40+10+200=250$. Above 250 the bonds are the fulcrum: at 255 they still receive 5. Below 250 the senior class becomes the fulcrum and the bonds receive nothing.
\item \emph{Without the priming assumption}, a DIP with only a junior lien on encumbered assets (sec.~364(c)(3)) ranks behind the senior lenders' collateral: at 220 the senior class would be paid in full and the DIP would receive 20 of its 40.
\item \emph{Loan-to-own at 300}: a plan pays the DIP and administrative claims from an exit facility of 50, reinstates the 200 of senior debt and gives the 50 of new equity to the bonds, 33.3 percent of face. An investor that bought 100 of the 150 bonds at 20 owns two-thirds of the new equity (33.3), a 66.7 percent gain.
\item \emph{Votes}: voting no, the investor stops the bond class accepting, since the rest hold only one-third of the amount. Voting yes, it carries the class only if more than half of the voting holders by number also accept; against ten small holders voting no, the class rejects.
\item \emph{Confirmation}: reinstated senior debt is unimpaired (sec.~1124(2)) and cannot supply the accepting impaired class, so a rejection by the bonds blocks this plan. A plan that changes the senior terms impairs the senior class; if it accepts, the plan can be confirmed over the bonds' rejection by cramdown, provided old equity receives nothing and the bonds get at least their liquidation value.
\end{itemize}
The investment is a bet on the valuation and on who controls the plan.

\qffield{Conditions} Performance statistics for distressed funds (US, 1990--2008: pooled net IRR 11.7 percent, multiple 1.56) are from Cambridge Associates as reported by Demaria and co-authors (p.~130). Their footnote on Europe and the rest of the world repeats, word for word, the figures given earlier for buyouts (p.~123), so those regional figures are not used. US law is cited from the US Code as published by the Legal Information Institute; case law on cramdown, the new-value exception, unfair discrimination and priming is not covered, and the waterfall's priority depends on the stated assumptions.

\qffield{Pitfalls} Buying a class that turns out to be out of the money. Equating a blocking position in a class vote with a veto over confirmation. Assuming a DIP loan primes existing liens without a sec.~364(d) order or consent. Underestimating the time and cost of a contested case. Treating distressed debt's low reported volatility as low risk.

\qffield{See also} \hyperref[def:private-debt-features]{private debt}; \hyperref[meth:private-capital-allocation]{allocation}; \hyperref[meth:restructuring-evaluation]{restructuring}.

\end{qfmethod}
